% =============================================================
%  ch29_machine_learning.tex  --  Klassisches ML & Deep Learning in ROS 2
% =============================================================

\chapter{Maschinelles Lernen in ROS\,2-Projekten}
\label{ch:machine_learning}

Kapitel~\ref{ch:merkmalserkennung} hat gezeigt, wie weit man in der
Bildverarbeitung \emph{ohne} Training kommt. Dieses Kapitel ergänzt die andere
Seite: \textbf{wann} sich maschinelles Lernen (ML) in einem ROS\,2-Roboter
lohnt, \textbf{wie} man Trainings- und Inferenzcode sauber von den ROS-Knoten
trennt und \textbf{womit} -- scikit-learn, SciPy, TensorFlow, PyTorch -- die
gängigen Verfahren (Regression, Clustering, Bayes, Autoencoder, ANN, CNN, PINN)
konkret angebunden werden. Der Anspruch ist nicht ein ML-Lehrbuch, sondern eine
\emph{Entscheidungs- und Integrationshilfe} für genau diese Roboterprojekte.

% -----------------------------------------------------------------
\section{Wann ML -- und wann nicht}
% -----------------------------------------------------------------
ML ist kein Selbstzweck. Die teuerste Fehlentscheidung ist, ein neuronales Netz
auf ein Problem zu werfen, das ein Kalman-Filter (Kap.~\ref{ch:merkmalserkennung}
hinweg zu Kap.~2) oder eine geschlossene Formel exakter, schneller und
nachvollziehbarer löst. Die nüchterne Faustregel:

\begin{infobox}[title={Faustregel: klassisch vor lernend}]
Gibt es ein \emph{physikalisches Modell} oder eine \emph{geometrische
Beziehung} (Kinematik, Projektion, Filtergleichung), nutze sie. Greife zu ML
erst, wenn (a) das Modell unbekannt oder zu komplex ist, (b) genügend
\emph{repräsentative Daten} vorliegen und (c) du den Inferenzaufwand im
Echtzeitbudget unterbringst.
\end{infobox}

\noindent Konkrete Anwendungsfälle in den Projekten dieses Buchs:

\begin{longtable}{@{}p{3.6cm}p{4.3cm}p{6.6cm}@{}}
\toprule
\textbf{Aufgabe} & \textbf{Klassisch reicht} & \textbf{ML lohnt sich, weil \ldots} \\
\midrule
Ball-/Objekt\-erkennung
  & Hintergrundsubtraktion, Farbschwelle, ORB
  & wechselnde Beleuchtung, Verdeckung, viele Objektklassen $\rightarrow$ \textbf{CNN}-Detektor \\
Flugbahn des Balls
  & ballistisches Modell + EKF
  & Magnus-Effekt/Spin schwer zu modellieren $\rightarrow$ \textbf{Regression/PINN} als Korrektur \\
Inverse Kinematik
  & analytisch/numerisch (Kap.~5)
  & redundante/elastische Mechanik, schnelle Näherung $\rightarrow$ \textbf{ANN} \\
Punktwolken-Seg\-mentierung
  & RANSAC-Ebene, Schwellwert
  & unstrukturierte Szene, mehrere Cluster $\rightarrow$ \textbf{Clustering} (DBSCAN) \\
Zustands\-/Anomalie\-überwachung
  & feste Schwellen
  & ``normal'' ist mehrdimensional und schwer zu fixieren $\rightarrow$ \textbf{Autoencoder} \\
Sensor\-kalibrierung
  & Least-Squares (SciPy)
  & nichtlineare, unbekannte Kennlinie mit Unsicherheit $\rightarrow$ \textbf{Gauß-Prozess} \\
\bottomrule
\end{longtable}

% -----------------------------------------------------------------
\section{Architektur: Training offline, Inferenz im Knoten}
% -----------------------------------------------------------------
Die wichtigste Designentscheidung ist die \textbf{Trennung von Training und
Inferenz}. Training ist datenintensiv, läuft batchweise auf einem Entwickler-PC
(gern mit GPU) und produziert eine \emph{Artefaktdatei} (Gewichte). Der
ROS\,2-Knoten lädt dieses Artefakt \emph{einmal} beim Hochfahren und führt im
Callback nur noch günstige Vorwärts-Inferenz aus.

\begin{center}
\begin{tikzpicture}[node distance=7mm and 16mm]
\node[node] (data) {Datensatz\\\scriptsize (rosbag2, Bilder, CSV)};
\node[node,right=of data] (train) {Training\\\scriptsize PyTorch/TF/sklearn (PC, GPU)};
\node[node,right=of train] (art) {Artefakt\\\scriptsize .pt / .onnx / .pkl};
\node[node,below=12mm of art] (node) {ML-Knote\\\code{ball\_detector}};
\node[topic,left=of node] (in) {/camera/image};
\node[topic,below=8mm of node] (out) {/ball/pose};
\draw[flow] (data) -- (train);
\draw[flow] (train) -- (art);
\draw[flow] (art) -- node[right,font=\scriptsize]{load() bei Start} (node);
\draw[flow] (in) -- node[above,font=\scriptsize]{Subscriber} (node);
\draw[flow] (node) -- node[right,font=\scriptsize]{Publisher} (out);
\draw[flow,dashed,accent2] (out) to[bend left=30]
  node[below,font=\scriptsize\itshape]{neue Daten sammeln} (data);
\end{tikzpicture}
\captionof{figure}{Offline-Training erzeugt ein Artefakt; der ROS\,2-Knoten lädt
es einmalig und betreibt nur noch Inferenz. Gesammelte Echtdaten fließen ins
nächste Training zurück.}
\end{center}

\begin{warnbox}[title={Inferenz nie im Spin-Thread blockieren}]
Ein \code{model(x)}-Aufruf kann zehn bis hunderte Millisekunden dauern. Läuft er
direkt im Subscriber-Callback des Standard-Executors, blockiert er alle anderen
Callbacks. Lege Inferenz in eine \textbf{eigene Callback-Group}
(\code{MutuallyExclusiveCallbackGroup}) mit \code{MultiThreadedExecutor} oder in
einen separaten Worker-Thread und veröffentliche das Ergebnis asynchron
(vgl.\ Kap.~\ref{ch:merkmalserkennung} zum Echtzeit-Scheduling).
\end{warnbox}

\begin{lstlisting}[style=py,caption={Grundgerüst eines Inferenz-Knotens (framework-agnostisch)}]
import rclpy
from rclpy.node import Node
from rclpy.callback_groups import MutuallyExclusiveCallbackGroup
from sensor_msgs.msg import Image
from geometry_msgs.msg import PointStamped
from cv_bridge import CvBridge

class BallDetector(Node):
    def __init__(self):
        super().__init__('ball_detector')
        # Modell EINMALIG laden (Pfad als ROS-Parameter):
        self.declare_parameter('model_path', 'model.onnx')
        path = self.get_parameter('model_path').value
        self.model = self.load_model(path)   # framework-spezifisch, s.u.
        self.bridge = CvBridge()
        cb = MutuallyExclusiveCallbackGroup()
        self.sub = self.create_subscription(
            Image, '/camera/image', self.on_image, 10, callback_group=cb)
        self.pub = self.create_publisher(PointStamped, '/ball/pose', 10)

    def on_image(self, msg):
        img = self.bridge.imgmsg_to_cv2(msg, 'rgb8')
        u, v, conf = self.infer(self.model, img)   # reine Vorwaertsinferenz
        if conf < 0.5:
            return
        out = PointStamped()
        out.header = msg.header                     # gleiche Zeit/Frame!
        out.point.x, out.point.y = float(u), float(v)
        self.pub.publish(out)
\end{lstlisting}

\noindent Drei Regeln, die in jedem ML-Knoten gelten sollten:
\begin{itemize}
\item \textbf{Modellpfad als Parameter}, nie hartkodiert -- so wechselt man
  Gewichte ohne Neukompilieren.
\item \textbf{Header übernehmen} (\code{stamp}, \code{frame\_id}): nur so passen
  Detektionen zeitlich zu \code{/tf} und lassen sich an einen EKF füttern.
\item \textbf{Konfidenz mitschätzen} und schwache Detektionen verwerfen statt zu
  raten.
\end{itemize}

% -----------------------------------------------------------------
\section{Die vier Bibliotheken und ihre Rolle}
% -----------------------------------------------------------------
\begin{longtable}{@{}p{2.7cm}p{4.0cm}p{7.8cm}@{}}
\toprule
\textbf{Bibliothek} & \textbf{Stärke} & \textbf{Rolle im ROS\,2-Projekt} \\
\midrule
\textbf{SciPy} & numerische Mathematik
  & Optimierung (\code{least\_squares}, \code{curve\_fit}), Filter
    (\code{signal}), Interpolation, lineare Algebra -- die Basis unter allem
    anderen, oft schon ausreichend. \\
\textbf{scikit-learn} & klassisches ML, ``flach''
  & Regression, Clustering, PCA, SVM, Gauß-Prozesse, Vorverarbeitung
    (\code{StandardScaler}, \code{Pipeline}). CPU, Modelle als
    \code{.pkl}/\code{joblib}. Erste Wahl bei tabellarischen/kleinen Daten. \\
\textbf{PyTorch} & Deep Learning, Forschung
  & ANN, CNN, Autoencoder, PINN. Dynamischer Graph, \code{torch.autograd}
    (wichtig für PINN), Export nach \code{.pt} oder ONNX. \\
\textbf{TensorFlow} & Deep Learning, Deployment
  & wie PyTorch; stark in \code{tf.keras} und Edge-Deployment
    (TF-Lite auf Raspberry Pi/Jetson). \\
\bottomrule
\end{longtable}

\begin{infobox}[title={Installation \& Workspace}]
Die ML-Pakete gehören \emph{nicht} in den ROS-Sourcetree. Üblich:
ein Python-\code{venv} oder \code{conda}-Env, in dem sowohl die ROS-Python-API
(über \code{rclpy}) als auch \code{torch}/\code{tensorflow} sichtbar sind.
In \code{package.xml} deklariert man die Abhängigkeiten via
\code{<exec\_depend>python3-scikit-learn</exec\_depend>} bzw.\ \code{rosdep};
große Frameworks (Torch/TF) installiert man per \code{pip} im Env und hält sie
aus \code{rosdep} heraus. \textbf{ONNX Runtime} ist der pragmatische
gemeinsame Nenner: in PyTorch/TF trainieren, nach \code{.onnx} exportieren, im
Knoten nur \code{onnxruntime} laden -- so muss der Roboter weder Torch noch TF
mitschleppen.
\end{infobox}

% -----------------------------------------------------------------
\section{Klassisches ML mit SciPy \& scikit-learn}
% -----------------------------------------------------------------

\subsection{Regression -- Kennlinien und Vorhersage}
Regression schätzt eine kontinuierliche Größe. Zwei typische Einsätze:
nichtlineare \emph{Sensorkalibrierung} mit SciPy und eine lernende
\emph{Trajektorienkorrektur} mit scikit-learn.

\begin{lstlisting}[style=py,caption={Nichtlineare Kalibrierkurve mit SciPy curve\_fit}]
import numpy as np
from scipy.optimize import curve_fit

# Modell: Distanz d aus Rohwert r eines IR-Sensors (1/r-Kennlinie)
def model(r, a, b, c):
    return a / (r + b) + c

popt, pcov = curve_fit(model, r_raw, d_true, p0=[1.0, 0.0, 0.0])
# popt -> in YAML speichern, im Knoten als feste Parameter laden:
d_pred = model(r_live, *popt)
\end{lstlisting}

\begin{lstlisting}[style=py,caption={Polynomiale Regression (sklearn-Pipeline)}]
from sklearn.pipeline import make_pipeline
from sklearn.preprocessing import PolynomialFeatures, StandardScaler
from sklearn.linear_model import Ridge
import joblib

# X: [Spin, Anstellwinkel, v0]  ->  y: Auftreffpunkt auf dem Tisch
reg = make_pipeline(StandardScaler(),
                    PolynomialFeatures(degree=2),
                    Ridge(alpha=1.0))
reg.fit(X_train, y_train)
joblib.dump(reg, 'landing_model.pkl')      # Artefakt fuer den Knoten
\end{lstlisting}

\subsection{Clustering -- Punktwolken und Objekte gruppieren}
Clustering ist \emph{unüberwacht}: keine Labels nötig. \code{DBSCAN} ist für
Robotik oft besser als \code{KMeans}, weil es die Clusterzahl nicht vorab
braucht und Rauschen explizit als Ausreißer markiert -- ideal, um aus einer
Tiefenpunktwolke Objekte zu trennen.

\begin{lstlisting}[style=py,caption={DBSCAN-Segmentierung einer Punktwolke}]
from sklearn.cluster import DBSCAN
import numpy as np

# pts: Nx3-Array aus sensor_msgs/PointCloud2
labels = DBSCAN(eps=0.05, min_samples=20).fit_predict(pts)
for k in set(labels):
    if k == -1:            # -1 = Rauschen
        continue
    cluster = pts[labels == k]
    centroid = cluster.mean(axis=0)   # -> als Objektkandidat publizieren
\end{lstlisting}

\subsection{Bayes \& Gauß-Prozesse -- Schätzen mit Unsicherheit}
Der bayessche Blick liefert nicht nur einen Wert, sondern eine
\emph{Verteilung}. Genau das passt zu Robotik, wo der EKF (Kap.~2) bereits
bayessch denkt. Ein \textbf{Gauß-Prozess} (GP) ist die elegante Form der
Regression \emph{mit} Konfidenzband: dort, wo wenig Daten liegen, wächst die
Unsicherheit -- ein direkt nutzbares Signal, um vorsichtig zu fahren.

\begin{lstlisting}[style=py,caption={Gauß-Prozess-Regression mit Unsicherheit (sklearn)}]
from sklearn.gaussian_process import GaussianProcessRegressor
from sklearn.gaussian_process.kernels import RBF, WhiteKernel

kernel = RBF(length_scale=1.0) + WhiteKernel(noise_level=1e-2)
gp = GaussianProcessRegressor(kernel=kernel, normalize_y=True)
gp.fit(X_train, y_train)

mu, sigma = gp.predict(X_query, return_std=True)
# sigma -> als Messrausch-Kovarianz direkt an den EKF weiterreichen
\end{lstlisting}

\noindent Der \textbf{naive Bayes}-Klassifikator (\code{sklearn.naive\_bayes})
bleibt die schnellste Baseline für diskrete Klassifikation (z.\,B.\ Zustand
``Greifer voll/leer'' aus wenigen Merkmalen) und eignet sich als
Vergleichsmaßstab, bevor man zu neuronalen Netzen greift.

% -----------------------------------------------------------------
\section{Neuronale Netze mit PyTorch}
% -----------------------------------------------------------------

\subsection{ANN (MLP) -- gelernte inverse Kinematik}
Ein mehrschichtiges Perzeptron (Multi-Layer Perceptron, MLP) approximiert
beliebige glatte Funktionen. Sinnvoll, wenn die analytische inverse Kinematik
mehrdeutig oder die reale Mechanik elastisch ist: das Netz liefert in
konstanter Zeit eine gute Startnäherung, die ein numerischer Schritt verfeinert.

\begin{lstlisting}[style=py,caption={MLP fuer inverse Kinematik (PyTorch)}]
import torch, torch.nn as nn

class IKNet(nn.Module):
    def __init__(self, n_in=3, n_out=4):   # (x,y,z) -> 4 Gelenkwinkel
        super().__init__()
        self.net = nn.Sequential(
            nn.Linear(n_in, 128), nn.ReLU(),
            nn.Linear(128, 128),  nn.ReLU(),
            nn.Linear(128, n_out))
    def forward(self, x):
        return self.net(x)

model = IKNet()
opt = torch.optim.Adam(model.parameters(), lr=1e-3)
loss_fn = nn.MSELoss()
for xb, yb in loader:                       # Training (offline, PC)
    opt.zero_grad()
    loss = loss_fn(model(xb), yb)
    loss.backward(); opt.step()
torch.save(model.state_dict(), 'ik_net.pt') # Artefakt
\end{lstlisting}

\subsection{CNN -- Objekt- und Ball-Erkennung}
Faltungsnetze (Convolutional Neural Networks) sind der Standard, sobald aus
\emph{Bildern} gelernt wird. Praktisch trainiert man selten von Null, sondern
nutzt \textbf{Transfer Learning}: ein vortrainiertes Backbone (z.\,B.\ ein
schlankes YOLO/MobileNet) wird auf die eigenen wenigen Objektklassen
feinjustiert. Im Knoten zählt nur die Inferenzlatenz -- daher Export nach ONNX
und Ausführung über \code{onnxruntime} (CPU/GPU/TensorRT).

\begin{lstlisting}[style=py,caption={CNN-Inferenz im Knoten via ONNX Runtime}]
import onnxruntime as ort
import numpy as np

class CNNDetector:
    def __init__(self, path):
        self.sess = ort.InferenceSession(
            path, providers=['CUDAExecutionProvider', 'CPUExecutionProvider'])
        self.inp = self.sess.get_inputs()[0].name

    def infer(self, img_rgb):                # img: HxWx3, uint8
        x = img_rgb.astype(np.float32).transpose(2, 0, 1)[None] / 255.0
        boxes, scores, classes = self.sess.run(None, {self.inp: x})
        return boxes, scores, classes
\end{lstlisting}

\subsection{Autoencoder -- Anomalie- und Zustandsüberwachung}
Ein Autoencoder lernt, seine Eingabe komprimiert zu rekonstruieren. Trainiert
\emph{nur auf Normalbetrieb}, steigt der Rekonstruktionsfehler bei
Anomalien -- ein label-freier Wächter für Vibrationen, Stromaufnahme oder
Gelenkmoment, der Defekte meldet, bevor feste Schwellen anschlagen.

\begin{lstlisting}[style=py,caption={Autoencoder zur Anomalieerkennung (PyTorch)}]
import torch, torch.nn as nn

class AE(nn.Module):
    def __init__(self, d):
        super().__init__()
        self.enc = nn.Sequential(nn.Linear(d, 16), nn.ReLU(), nn.Linear(16, 4))
        self.dec = nn.Sequential(nn.Linear(4, 16), nn.ReLU(), nn.Linear(16, d))
    def forward(self, x):
        return self.dec(self.enc(x))

# Im Knoten: Fehler gegen Schwelle aus dem Training (z.B. 99%-Quantil)
def is_anomaly(model, x, thresh):
    with torch.no_grad():
        err = ((model(x) - x) ** 2).mean().item()
    return err > thresh, err
\end{lstlisting}

\subsection{PINN -- physik-informierte Netze}
Ein \textbf{Physics-Informed Neural Network} (PINN) verbindet beide Welten: Die
Verlustfunktion bestraft nicht nur die Datenabweichung, sondern auch die
Verletzung einer bekannten \emph{Differentialgleichung}. Für die Ballflugbahn
mit Luftwiderstand und Magnus-Effekt heißt das: Das Netz $\bm{x}(t)$ wird so
trainiert, dass es zu den Messpunkten passt \emph{und} die Bewegungsgleichung
$m\ddot{\bm{x}} = \bm{F}_g + \bm{F}_{\text{drag}} + \bm{F}_{\text{magnus}}$
erfüllt. Vorteil: physikalisch plausible Vorhersagen auch mit wenigen Daten.
Der Schlüssel ist \code{torch.autograd.grad}, das die Zeitableitungen
$\dot{\bm{x}}, \ddot{\bm{x}}$ exakt aus dem Netz liefert.

\begin{lstlisting}[style=py,caption={PINN-Residuum fuer eine Bewegungsgleichung (PyTorch)}]
import torch

def physics_residual(net, t, m, g):
    t = t.requires_grad_(True)
    x = net(t)                                   # Position(t), z.B. Hoehe
    v = torch.autograd.grad(x, t, torch.ones_like(x), create_graph=True)[0]
    a = torch.autograd.grad(v, t, torch.ones_like(v), create_graph=True)[0]
    # ODE-Residuum:  m*a + drag(v) - (-m*g) = 0  -> hier vereinfacht ohne drag
    return m * a + m * g

# Gesamtverlust = Daten-MSE + lambda * mean(residual^2)
loss = mse(net(t_data), x_data) + lam * (physics_residual(net, t_col, m, g)**2).mean()
\end{lstlisting}

% -----------------------------------------------------------------
\section{Verfahren-Wegweiser}
% -----------------------------------------------------------------
\begin{longtable}{@{}p{2.6cm}p{2.8cm}p{3.0cm}p{5.4cm}@{}}
\toprule
\textbf{Verfahren} & \textbf{Typ} & \textbf{Werkzeug} & \textbf{Wofür im Roboter} \\
\midrule
Least-Squares / \code{curve\_fit} & Fit & SciPy &
  Kalibrierung, Modellabgleich \\
Lineare/Ridge-Regression & überwacht & scikit-learn &
  einfache Vorhersage, Baseline \\
Gauß-Prozess & überwacht, bayessch & scikit-learn &
  Regression \emph{mit} Unsicherheit für den EKF \\
Naive Bayes / SVM & überwacht & scikit-learn &
  schnelle Klassifikation weniger Merkmale \\
KMeans / DBSCAN & unüberwacht & scikit-learn &
  Punktwolken-/Objektsegmentierung \\
PCA & unüberwacht & scikit-learn &
  Dimensionsreduktion, Merkmalsanalyse \\
MLP (ANN) & überwacht, tief & PyTorch/TF &
  gelernte IK, nichtlineare Regelung \\
CNN & überwacht, tief & PyTorch/TF + ONNX &
  Bild-/Objekt-/Ballerkennung \\
Autoencoder & unüberwacht, tief & PyTorch/TF &
  Anomalie-/Zustandsüberwachung \\
PINN & hybrid & PyTorch &
  Trajektorie mit Physik-Nebenbedingung \\
\bottomrule
\end{longtable}

\begin{infobox}[title={Kernbotschaft}]
Maschinelles Lernen ist in ROS\,2 kein Bruch, sondern ein weiterer Knoten:
\emph{offline lernen, ein Artefakt exportieren, im Lifecycle-Knoten laden,
Inferenz aus dem kritischen Pfad heraushalten und Ergebnisse mit korrektem
Header publizieren}. Beginne mit SciPy/scikit-learn; greife zu PyTorch/TF erst,
wenn Bilder, hohe Dimensionen oder Physik-Nebenbedingungen es erzwingen.
\end{infobox}
